\section{正项级数的比较}
\label{sec:02-Calculus/05-Sequences-and-Series/03}

\subsection{正项带来的特殊便利}

如果所有项 \(a_n\ge0\)，部分和 \(S_N\) 的行为就非常简单：每多加一项，和要么不变，要么往上走。部分和数列单调不减。

这个简单事实给出了一个强有力的等价条件：正项级数收敛，当且仅当它的部分和有上界。因为单调不减的数列，如果还能被某个数压住，就必然收敛到某个不超过该上界的极限；如果没有上界，它就直奔无穷。

这是所有正项级数判别法的公共出发点。每种方法都在做同一件事：找到某种方式证明部分和有限，或者证明它无限。

\subsection{需要一个锚：什么样的级数我们已经知道敛散？}

在谈论比较之前，你手里必须有一些已知收敛或发散的参照级数。目前最牢靠的锚有两个：

\begin{itemize}
\tightlist
\item 几何级数 \(\sum r^n\)：\(|r|<1\) 收敛，\(|r|\ge1\) 发散。部分和可以直接写出来验证。
\item 调和级数 \(\sum 1/n\)：发散。上一节的分组论证已经确认。
\end{itemize}

但这两个锚还不够——很多级数的衰减速度介于几何级数和调和级数之间。比如 \(\sum 1/n^2\)，它比调和级数衰减得快，但还没有指数级快。是收敛还是发散？几何级数和调和级数的知识回答不了这个问题。需要一个新工具来判断 \(p\neq1\) 的 \(p\) 级数。

\subsection{积分判别：用连续的面积夹逼离散的和}

对于形如 \(\sum f(n)\) 的正项级数，若 \(f(x)\) 在 \([1,\infty)\) 上为正、连续且递减，那么函数曲线下方的面积就夹在矩形面积之间，可以用来判断级数的敛散。

具体说，比较 \(f\) 在 \([1,N]\) 上的积分与矩形面积：
\[
\int_1^N f(x)\,dx
\le \sum_{n=1}^{N-1}f(n)
\le f(1)+\int_1^{N-1}f(x)\,dx.
\]

直观理解就是：左端矩形（高为各左端点函数值）的面积和大于曲线下方面积；右端矩形（高为各右端点函数值）的面积和小于曲线下方面积。让 \(N\to\infty\)，级数与反常积分的敛散性绑定在一起。

立即检验关键例 \(\sum 1/n^p\)。令 \(f(x)=x^{-p}\)，积分
\[
\int_1^\infty\frac{dx}{x^p}
\]
当且仅当 \(p>1\) 收敛。所以 \(\sum 1/n^p\) 也当且仅当 \(p>1\) 收敛。

至此我们有了一个系统性的 \(p\)-级数锚：\(p>1\) 收敛，\(p\le1\) 发散。\(\sum 1/n^2\) 收敛，\(\sum 1/\sqrt{n}\) 发散。边界 \(p=1\) 落在发散一侧。

积分判别还有一个附带好处：余项可以给出不等式估计：
\[
\int_{N+1}^{\infty}f(x)\,dx
\le R_N=\sum_{n=N+1}^{\infty}f(n)
\le\int_N^{\infty}f(x)\,dx.
\]

如果收敛，这个不等式给了你截断误差的上下界。

\subsection{直接比较：夹在两已知级数之间}

有了 \(p\)-级数和几何级数这两种锚，就可以用大小关系判断其他级数了。

直接比较的规则很简单：若从某项开始恒有 \(0\le a_n\le b_n\)，则 \(\sum b_n\) 收敛推出 \(\sum a_n\) 收敛；反过来，若 \(a_n\ge b_n\ge0\) 且 \(\sum b_n\) 发散，则 \(\sum a_n\) 发散。

例如：
\[
0<\frac1{n^2+1}<\frac1{n^2},
\]
\(\sum 1/n^2\) 收敛，所以 \(\sum 1/(n^2+1)\) 收敛。

但方向必须对。用一个比 \(a_n\) 大却发散的级数压在它上面说明不了什么——比如 \(\sum 1/n\) 发散，用它去压 \(\sum 1/n^2\)（\(1/n^2<1/n\)）不能得出结论。同理，用一个比 \(a_n\) 小却收敛的级数垫在下面也说明不了什么。不等式方向必须和结论方向一致：大者收敛则小者收敛；小者发散则大者发散。

\subsection{极限比较：当直接不等式不好列}

有时候 \(a_n\) 和某个参照级数的 \(b_n\) 没有整齐的大小关系，但它们的比值趋于一个正有限常数。这时可以绕过构造不等式的麻烦。

若
\[
\lim_{n\to\infty}\frac{a_n}{b_n}=c,\qquad 0<c<\infty,
\]
则 \(\sum a_n\) 与 \(\sum b_n\) 同时收敛或同时发散。

背后道理：当比值趋于 \(c>0\)，对充分大的 \(n\) 有 \(c/2 < a_n/b_n < 2c\)，即 \(a_n\) 被夹在 \(b_n\) 的常数倍之间。乘以常数不改变级数的敛散性，所以两者命运相同。正有限比值说明尾部是同一数量级。

例如 \(a_n = (n+1)/(n^3+2)\)。当 \(n\) 很大时，分子的 \(n\) 和分母的 \(n^3\) 主导，所以 \(a_n \sim n/n^3 = 1/n^2\)。精确验证：
\[
\frac{a_n}{1/n^2}=\frac{n^2(n+1)}{n^3+2}=\frac{n^3+n^2}{n^3+2}\to1,
\]
正的有限数。\(\sum 1/n^2\) 收敛，所以原级数收敛。

极限比较让判断从``找不等式''变成``找主导项''——后者在很多场合更容易直接读出。

\subsection{比值与根值：检测指数尺度的衰减}

积分判别和比较判别适合判断代数衰减（\(n^{-p}\) 型）。但当项里出现阶乘、指数和 \(n\) 次幂时，衰减速度可能远远超过任何代数幂次——这时需要直接测量相邻项的乘法比例。

比值检验：若
\[
L=\lim_{n\to\infty}\left|\frac{a_{n+1}}{a_n}\right|,
\]
则 \(L<1\) 时级数绝对收敛，\(L>1\) 时发散。背后的比较对象是几何级数：\(L<1\) 意味着尾部项的衰减速度最终快于某个公比小于 1 的几何级数。

例如 \(a_n = n!\,/\,n^n\)：
\[
\frac{a_{n+1}}{a_n}
=\frac{(n+1)!}{(n+1)^{n+1}}\cdot\frac{n^n}{n!}
=\frac{n+1}{(n+1)^{n+1}}\cdot n^n
=\frac{n^n}{(n+1)^n}
=\frac1{(1+1/n)^n}\to\frac1e<1,
\]
级数收敛。注意项的衰减比任何几何级数都快，比值检验干净地捕捉到了这一点。

根值检验是同样的逻辑，但从另一个角度：若
\[
L=\limsup_{n\to\infty}|a_n|^{1/n},
\]
则 \(L<1\) 收敛、\(L>1\) 发散。根值检验把第 \(n\) 项与 \(L^n\) 比较——同样是与几何级数的间接对话。\(\limsup\) 是``尾部上确界的极限''，总存在；当普通极限存在时两者相等，可以先按普通极限理解。

\subsection{\(L=1\) 不是答案，是工具的极限}

比值检验和根值检验都有一条免责声明：当 \(L=1\) 时，没有结论。这不是级数的问题，是工具的分辨率不足。

\(\sum 1/n\) 发散，比值极限为
\[
\frac{1/(n+1)}{1/n}=\frac{n}{n+1}\to1.
\]
\(\sum 1/n^2\) 收敛，比值极限同样是
\[
\frac{1/(n+1)^2}{1/n^2}=\frac{n^2}{(n+1)^2}\to1.
\]

两个级数穷尽了两种可能——一个发散一个收敛——比值检验却对它们都说``我不知道''。这是因为比值和根值检验是把级数与几何级数做比较，而任何代数衰减的项都会产生 \(L=1\)。\(L=1\) 意味着该级数的衰减或增长既不是指数级的也不是超指数级的，而这种情况下你需要回过头求助于更精细的比较或积分判别。

\subsection{判据选择：看尾部什么结构}

面对一个具体的正项级数，选择判别法可以参考下面的线索：
\begin{itemize}
\tightlist
\item 项与 \(1/n^p\) 或简单有理式同阶：用比较判别或极限比较判别，参照已知的 \(p\)-级数。
\item 项来自某个正的递减函数：积分判别既能判断敛散，还能附带余项估计。
\item 项含阶乘、指数、\(n\) 的 \(n\) 次幂：比值或根值检验，它们能探测超几何级别的衰减。
\item 做任何分析之前：先扫一眼通项是否趋于零。如果不趋于零，直接发散，不需要浪费其他判据的时间。
\end{itemize}

这四种方法是同一个问题的不同切面——尾部有多重？几何衰减、代数衰减、对数衰减对应不同的工具灵敏度。选择工具的本质是选择你相信的标杆：几何级数还是 \(p\)-级数。

\subsection{收敛不等于能算出和}

最后这个提醒容易在技巧的兴奋中被忽略：所有判别法只回答``收敛还是发散''，不负责算出极限值。\(\sum 1/n^2=\pi^2/6\)，这个结果不是通过比较判别或积分判别得到的——它需要一个完全不同的论证。很多收敛级数的和都没有初等闭式。判别法告诉你这条路走得通，但走到哪里，有时候连数学本身都只能给数值近似。
